\section{Goals and Role of the Prototype}

Zweck des Prototyps als Forschungsplattform, nicht als fertiges Produkt, kontrollierbare Untersuchungsumgebung für algorithmische Textursynthese, reproduzierbare Experimente, Vergleich verschiedener Ansätze unter identischen Bedingungen, Unterstützung der Forschungsfragen, Verbindung zwischen Theorie und Evaluation, Nutzen für Forschung, Game Development, Procedural Content Generation, Echtzeitgrafik, mögliche Nutzung durch Studierende, Forschende oder Entwickler für Experimente mit Texturalgorithmen

\section{Analysis of Existing Tools and Frameworks}

(kurzes Kapitel)

existierende Tools und Engines für prozedurale Texturen, Substance Designer, Blender Shader Nodes, Unreal Engine Materials, Houdini, bestehende Research-Prototypes, Vorteile existierender Lösungen wie ausgereifte Node-Systeme oder GPU-Integration, Einschränkungen bezüglich Kontrollierbarkeit, Vergleichbarkeit, reproduzierbarer Messungen, fehlender Fokus auf kombinierte Untersuchung von procedural/rule-based/optimization-based Ansätzen, fehlende direkte Integration der gewünschten Evaluationsmetriken, Begründung für eigene spezialisierte Implementation angepasst auf Forschungsfragen und Experimentdesign

\section{System Architecture}

\subsection{Technology Stack and Framework Selection}

Next.js, React, React Three Fiber, Three.js, TypeScript, WebGL, Browser-basierte Architektur, Echtzeitdarstellung im Browser, schnelle Iteration, Visualisierung und Interaktivität, einfache Parametrisierung, Plattformunabhängigkeit, Fokus auf Forschungsprototyp statt produktionsreifer Engine

\subsection{High-Level System Architecture}

Aufbau der Plattform, Rendering-Pipeline, Terrain-System, Objekt-System, Texture-Generation-Pipeline, UI- und Parameter-System, Datenfluss zwischen Parametern, Texturerzeugung und Rendering, modulare Architektur zur Erweiterbarkeit weiterer Syntheseverfahren und Evaluationsmethoden

\section{Procedural Environment Foundation}

\subsection{Terrain Generation}

Noise-basierte Terrain-Erzeugung, fBm, Ridge Noise, Turbulence, Kombination mehrerer Noise-Layer, Chunk-System, deterministische Generierung mittels Seeds, Ziel einer variablen aber kontrollierbaren Umgebung

\subsection{Object Placement and Scene Composition}

deterministische Objektplatzierung, verschiedene Objekttypen, kontrollierte Variation, reproduzierbare Szenen, Material-Slots, Nutzung als Testumgebung für unterschiedliche Texturtypen und Parameterkombinationen

\subsection{Real-Time Rendering Pipeline}

Mesh-Erzeugung, Materialsystem, Texture-Canvas-Generierung, Rendering-Ablauf, Echtzeitaktualisierung von Parametern, Browser-basierte Darstellung, Zusammenhang zwischen Texturerzeugung und Renderkosten

\section{Texture Synthesis Implementation}

\subsection{Implemented Procedural Methods}

Perlin/fBm, Domain Warping, Worley/Voronoi, mathematische Grundlagen kurz implementierungsbezogen, Parameterisierung der Algorithmen, Einfluss von Scale, Frequency, Variance, Seeds und Layering, Unterschiede visueller Eigenschaften und Strukturen

\subsection{Rule-Based and Parameterized Texture Models}

regelbasierte Definition von Texturverhalten, Material-Slots, Kombination von Regeln und Noise-Funktionen, Parametrisierung von Farben und Strukturen, kontrollierbare Variationen, Reproduzierbarkeit, Ziel einer flexiblen Texture-Pipeline

\subsection{Texture Approximation and Optimization Concepts}

Ansatz zur Approximation von Zieltexturen, Parametersuche, Vergleich generierter und gewünschter Texturen, mögliche Optimierungsstrategien, Einschränkungen der aktuellen Arbeit, Fokus auf kontrollierbare Parameter statt vollständige Reproduktion realer Texturen

\section{Performance and Scalability Considerations}

\subsection{Runtime Challenges}

CPU-Last durch Texture-Generierung, Chunk-Loading-Probleme, Render-Lag, Speicherverbrauch, Echtzeit-Anforderungen, Browser-Limitationen, Trade-offs zwischen Qualität und Performance

\subsection{Optimization Strategies}

Caching von Geometrie und Materialien, Texture-Reuse-Rate, React.memo, Chunk-Streaming, Preloading, reduzierte Segmentanzahl, Performanceverbesserungen, Einfluss auf Skalierbarkeit und Reaktionsfähigkeit des Systems

\subsection{Known Limitations and Technical Constraints}

Worley-Artefakte, Main-Thread-Limitierungen, fehlende GPU-Compute-Shaders, Browserabhängigkeit, eingeschränkte Parallelisierung, Fokus der Arbeit auf Forschungsziele statt vollständige Produktionsoptimierung

\section{Experimental and Evaluation Infrastructure}

Before conducting the final evaluation, the prototype was inspected to verify which experimental data could be collected reliably. This step was necessary because the implemented system contains several different texture synthesis approaches with different technical characteristics. Some methods are implemented as GPU-based WebGPU shaders, while others are executed on the CPU or in a web worker. As a result, not all methods support the same resolutions, runtime behavior, or reproducibility guarantees.

The inspection showed that the implemented methods expose different parameter sets and export options. For example, the optimization-based method can export parameter configurations and diagnostic information, while the procedural methods are mainly evaluated through generated texture images, parameter settings, and runtime measurements.

This verification step helped define the final evaluation setup. A common resolution of 256x256 pixels was selected for comparison because this resolution is supported by all implemented methods. Higher resolutions are considered only for methods that technically support them. The optimization-based approach currently focuses on the Wood texture model. Therefore, optimization results are interpreted as an investigation of parameter estimation for one representative procedural material model rather than as a general optimization framework for all implemented synthesis methods.

\subsection{Measurement and Monitoring System}

The prototype provides several mechanisms for collecting evaluation data. Runtime-related values such as texture generation time and rendering performance can be measured using browser-based timing functions. In addition, the terrain environment includes monitoring functionality for values such as frames per second, texture generation time, chunk generation time, and render time.

The collected measurement data is intended to support comparison between the implemented synthesis methods under controlled conditions. Since GPU-based and CPU-based methods have different execution characteristics, the evaluation records the execution context of each method separately. This distinction is important because WebGPU shader methods primarily depend on GPU execution, while example-based methods such as pixel-based synthesis or patch-based synthesis depend more strongly on CPU performance and memory access patterns.

\subsection{Evaluation Metrics and Comparison Criteria}

The optimization method uses a weighted fitness score based on several image statistics. These include mean brightness, variance, gradient ratio, edge density, and histogram similarity. The individual components are combined into a single scalar fitness value, where higher values indicate a closer match between the generated texture and the target texture.

This metric was selected because exact pixel-level reproduction is not realistic for procedural texture synthesis. Procedural models often generate visually similar structures at different pixel positions. Therefore, statistical and structural characteristics are more suitable than a purely pixel-based error metric for comparing generated and target textures.

\subsection{Reproducibility and Controlled Experiments}

Reproducibility is an important requirement for the evaluation because the prototype is intended as a research platform rather than a production-ready texture tool. The implementation was inspected to determine which methods support deterministic generation. The procedural shader-based methods use seed-based generation and are therefore expected to reproduce the same output when the same parameters and seed are used. The optimization-based method also uses seeded random values, which makes repeated optimization runs reproducible under identical settings.

However, not all implemented methods provide the same level of reproducibility. The Efros-Leung implementation contains random choices that are not fully controlled by the configured seed. This means that repeated runs may produce slightly different results even when the visible parameters remain unchanged. This behavior is a limitation and is considered when interpreting the reproducibility results.

For the controlled experiments, seeds, texture resolution, target images, and parameter settings are kept constant whenever possible. This allows the influence of individual methods and parameters to be analyzed more clearly.

\section{Summary of the Implementation}

Zusammenfassung der entwickelten Forschungsplattform, wichtigste technische Erkenntnisse, Nutzen der Plattform für die Untersuchung algorithmischer Textursynthese, Beitrag zur Beantwortung der Forschungsfragen, Übergang zur Evaluation und Analyse der Ergebnisse